% PRISMA 2020-style full demonstration template
% Edit only the data and labels in the CONFIGURATION section.
% All example numbers are illustrative, not findings from a real review.
\documentclass[border=4pt]{standalone}
\usepackage[T1]{fontenc}
\usepackage{lmodern}
\usepackage{tikz}
\usetikzlibrary{arrows.meta,calc,positioning,fit}
\usepackage{xcolor}
\usepackage{tcolorbox}

% ==================== CONFIGURATION ====================
% Search sources (change counts and database names here)
\newcommand{\SourceOne}{Scopus}
\newcommand{\SourceTwo}{PubMed}
\newcommand{\SourceThree}{Web of Science}
\def\NSourceOne{142}
\def\NSourceTwo{38}
\def\NSourceThree{20}

% Records removed before screening
\def\NDuplicates{12}
\def\NAutomation{3}
\def\NOtherRemoved{5}

% Screening exclusions and unavailable full texts
\def\NScreenExcluded{150}
\def\NNotRetrieved{2}

% Report-level exclusion reasons (edit descriptions and counts)
\newcommand{\ReasonOne}{Ineligible study design}
\newcommand{\ReasonTwo}{Ineligible outcomes}
\newcommand{\ReasonThree}{Insufficient relevant data}
\def\NReasonOne{4}
\def\NReasonTwo{3}
\def\NReasonThree{1}

% Calculated totals: no manual recounting required
\pgfmathtruncatemacro{\NIdentified}{\NSourceOne+\NSourceTwo+\NSourceThree}
\pgfmathtruncatemacro{\NRemoved}{\NDuplicates+\NAutomation+\NOtherRemoved}
\pgfmathtruncatemacro{\NScreened}{\NIdentified-\NRemoved}
\pgfmathtruncatemacro{\NSought}{\NScreened-\NScreenExcluded}
\pgfmathtruncatemacro{\NAssessed}{\NSought-\NNotRetrieved}
\pgfmathtruncatemacro{\NReportsExcluded}{\NReasonOne+\NReasonTwo+\NReasonThree}
\pgfmathtruncatemacro{\NIncluded}{\NAssessed-\NReportsExcluded}

% Independent review/study counts may differ in real searches where
% multiple reports correspond to one study. Adapt labels accordingly.
\newcommand{\IncludedLabel}{Studies included in review}

% Theme
\definecolor{prismNavy}{HTML}{174A7E}
\definecolor{prismInk}{HTML}{202B36}
\definecolor{prismTeal}{HTML}{007C83}
\definecolor{prismOrange}{HTML}{B85C00}
\definecolor{prismOrangeLight}{HTML}{FFF0D9}
\definecolor{prismGreen}{HTML}{287A3D}
\definecolor{prismGreenLight}{HTML}{E1F2E5}

% ================= END CONFIGURATION ===================

\begin{document}
\tcbox[colback=white,colframe=prismNavy,arc=4pt,boxrule=1.2pt,
  boxsep=2pt,left=2pt,right=2pt,top=2pt,bottom=2pt]{%
\begin{tikzpicture}[
    font=\sffamily\normalsize,
    text=prismInk,
    basic/.style={rectangle,draw=prismNavy,very thick,rounded corners=3pt,
      align=left,text width=5.20cm,minimum height=1.20cm,inner sep=7pt,fill=white},
    exclusion/.style={basic,draw=prismOrange,fill=prismOrangeLight},
    outcome/.style={basic,draw=prismGreen,fill=prismGreenLight},
    phase/.style={rectangle,draw=prismTeal,fill=prismTeal,text=white,
      very thick,rounded corners=3pt,align=center,text width=2.4cm,
      minimum height=0.88cm,inner sep=5pt,font=\sffamily\bfseries},
    headline/.style={rectangle,draw=prismNavy,fill=prismNavy,text=white,
      very thick,rounded corners=3pt,align=center,text width=12.0cm,
      minimum height=0.75cm,inner sep=7pt,font=\sffamily\bfseries\large},
    link/.style={-{Latex[length=2.5mm]},line width=1.4pt,draw=prismNavy,
      shorten >=0pt,shorten <=0pt}
]
% Header, phase labels and two-column PRISMA flow
\node[headline] (heading) at (3.53,0.65)
 {Identification of studies via databases and registers};
\node[phase] (phaseID) at (-3.35,-1.65) {Identification};
\node[basic] (records) at (1.15,-1.65)
 {\textbf{Records identified from:}\\
 \SourceOne\ ($n=\NSourceOne$)\\
 \SourceTwo\ ($n=\NSourceTwo$)\\
 \SourceThree\ ($n=\NSourceThree$)\\
 \textbf{Total ($n=\NIdentified$)}};
\node[exclusion] (preremoved) at (8.1,-1.65)
 {\textbf{Records removed before\\screening:}\\
 Duplicate records ($n=\NDuplicates$)\\
 Automation exclusions ($n=\NAutomation$)\\
 Other reasons ($n=\NOtherRemoved$)\\
 \textbf{Total ($n=\NRemoved$)}};

\node[phase] (phaseScreen) at (-3.35,-5.10) {Screening};
\node[basic] (screen) at (1.15,-5.10)
 {\textbf{Records screened}\\($n=\NScreened$)};
\node[exclusion] (screenout) at (8.1,-5.10)
 {\textbf{Records excluded}\\($n=\NScreenExcluded$)};

\node[basic] (sought) at (1.15,-7.05)
 {\textbf{Reports sought for retrieval}\\($n=\NSought$)};
\node[exclusion] (unavailable) at (8.1,-7.05)
 {\textbf{Reports not retrieved}\\($n=\NNotRetrieved$)};

\node[basic] (assessed) at (1.15,-9.8)
 {\textbf{Reports assessed for eligibility}\\($n=\NAssessed$)};
\node[exclusion] (reasons) at (8.1,-9.8)
 {\textbf{Reports excluded:\\(with reasons):}\\
 \ReasonOne\ ($n=\NReasonOne$)\\
 \ReasonTwo\ ($n=\NReasonTwo$)\\
 \ReasonThree\ ($n=\NReasonThree$)\\
 \textbf{Total ($n=\NReportsExcluded$)}};

\node[phase] (phaseInclude) at (-3.35,-13.0) {Included};
\node[outcome] (final) at (1.15,-13.0)
 {\textbf{\IncludedLabel}\\($n=\NIncluded$)};

% Draw after nodes, with exact boundaries rather than fixed coordinates
\draw[link] (records.east) -- (preremoved.west);
\draw[link] (records.south) -- (screen.north);
\draw[link] (screen.east) -- (screenout.west);
\draw[link] (screen.south) -- (sought.north);
\draw[link] (sought.east) -- (unavailable.west);
\draw[link] (sought.south) -- (assessed.north);
\draw[link] (assessed.east) -- (reasons.west);
\draw[link] (assessed.south) -- (final.north);

% Outer frame is supplied by the original navy tcolorbox above.
\end{tikzpicture}%
}%
\end{document}